% ============================================================
%  PHỤ LỤC: BẢNG PHÂN CÔNG CÔNG VIỆC
% ============================================================
\chapter{BẢNG PHÂN CÔNG CÔNG VIỆC}
\label{appendix:phan_cong}

%
% REQUIREMENT (từ main.md – Bảng phân công):
%   - Liệt kê đầy đủ toàn bộ thành viên nhóm.
%   - Phân công nhiệm vụ theo từng giai đoạn:
%       + Phân tích yêu cầu
%       + Thiết kế hệ thống
%       + Lập trình (ghi rõ service/module nào)
%       + Kiểm thử
%       + Viết báo cáo
%   - Ghi rõ mức độ đóng góp (%) của từng thành viên.
%   - Có thể bổ sung cột: Nhận xét của nhóm trưởng.
%
%   Gợi ý format bảng:
%

\begin{longtable}{|c|p{3.2cm}|p{8.0cm}|c|}
  \hline
  \textbf{STT} & \textbf{Họ và tên / MSSV} & \textbf{Nhiệm vụ phân công chi tiết} & \textbf{Đóng góp} \\
  \hline
  \endfirsthead

  \hline
  \textbf{STT} & \textbf{Họ và tên / MSSV} & \textbf{Nhiệm vụ phân công chi tiết} & \textbf{Đóng góp} \\
  \hline
  \endhead

  \hline
  \endfoot

  1 & \textbf{Đặng Xuân Lâm}\newline N22DCCN047 & 
  \textbf{1. Phân tích \& Thiết kế:}
  \begin{itemize}[leftmargin=*, nosep, topsep=0pt]
      \item Thiết kế kiến trúc Microservices tổng thể và quy hoạch Database-per-service (PostgreSQL).
      \item Phân tích luồng giao tiếp bất đồng bộ: Event-Driven Architecture, Saga Orchestration, luồng đồng bộ dữ liệu (Data Pipeline).
      \item Thiết kế lược đồ CSDL cho hệ thống giao dịch cốt lõi (Booking, Payment).
  \end{itemize}
  \vspace{0.1cm}
  \textbf{2. Lập trình (Backend/Infra):}
  \begin{itemize}[leftmargin=*, nosep, topsep=0pt]
      \item Xây dựng \texttt{BookingService}: Cài đặt MassTransit State Machine điều phối vòng đời đặt phòng khắt khe, xử lý concurrency chống Double-booking.
      \item Xây dựng \texttt{PaymentService}: Tích hợp cổng thanh toán VNPay/Stripe, xử lý Webhook an toàn, đối soát giao dịch.
      \item Thiết lập \textbf{RabbitMQ} (Exchange/Queue) và triển khai \textbf{Transactional Outbox Pattern}.
      \item Cấu hình hạ tầng \textbf{Debezium \& Apache Kafka} phục vụ Change Data Capture (CDC).
      \item Container hóa toàn bộ hệ thống bằng \textbf{Docker Compose}.
  \end{itemize}
  \vspace{0.1cm}
  \textbf{3. Báo cáo \& Trình bày:}
  \begin{itemize}[leftmargin=*, nosep, topsep=0pt]
      \item Soạn thảo Báo cáo (Chương 2: Thiết kế hệ thống).
      \item Thiết lập \textbf{.NET Aspire Dashboard} phục vụ Distributed Tracing (OpenTelemetry) cho buổi Demo.
      \item Xây dựng kịch bản trình bày luồng Saga và xử lý bù trừ (Compensation).
  \end{itemize} & 50\% \\
  \hline
  2 & \textbf{Phan Nhật Minh}\newline N22DCCN054 & 
  \textbf{1. Phân tích \& Thiết kế:}
  \begin{itemize}[leftmargin=*, nosep, topsep=0pt]
      \item Phân rã Bounded Contexts, thiết kế hợp đồng API (API Contracts) cho hệ sinh thái.
      \item Thiết kế kiến trúc API Gateway làm điểm chạm duy nhất (Single Point of Entry).
      \item Thiết kế UX/UI cho ứng dụng Frontend (Admin Dashboard \& Guest Portal).
      \item Thiết kế cấu trúc dữ liệu NoSQL (Inverted Index) phục vụ tối ưu hóa tìm kiếm.
  \end{itemize}
  \vspace{0.1cm}
  \textbf{2. Lập trình (Backend/Frontend):}
  \begin{itemize}[leftmargin=*, nosep, topsep=0pt]
      \item Cài đặt \textbf{YARP API Gateway}: Định tuyến (Routing), xác thực (Auth offloading), Rate Limiting.
      \item Xây dựng \texttt{SearchService}: Tích hợp \textbf{Elasticsearch}, lập chỉ mục (Indexing) realtime từ Kafka, xử lý truy vấn tọa độ \& bộ lọc đa chiều.
      \item Xây dựng \texttt{PropertyService} (Quản lý phòng/thuế), \texttt{UserService} (Identity/JWT), và \texttt{ChatService} (Real-time message).
      \item Áp dụng triệt để kiến trúc \textbf{CQRS} (MediatR) và \textbf{FastEndpoints} cho mọi API.
      \item Hoàn thiện toàn bộ giao diện \textbf{Frontend (React, Zustand, TanStack Query, Shadcn/UI)}.
  \end{itemize}
  \vspace{0.1cm}
  \textbf{3. Báo cáo \& Trình bày:}
  \begin{itemize}[leftmargin=*, nosep, topsep=0pt]
      \item Viết Báo cáo (Chương 3: Triển khai \& Chương 4: Kết luận).
      \item Cấu trúc lại source code theo Clean Architecture, thiết kế Slide thuyết trình.
      \item Trình diễn giao diện người dùng và giải thích cấu trúc CQRS trong buổi bảo vệ.
  \end{itemize} & 50\% \\
  \hline
\end{longtable}
